%% The full discussion of each limitation, moved out of \Cref{sec:limits} to keep the main text
%% short.  Shared by ../paper/main.tex and ../submission/body.tex.
We list the boundaries of the result plainly, because several of them are load-bearing.

\paragraph{Memory is $O(d^2)$, not $O(d)$.}
The $O(dN)$ claim is about \emph{time}. The method stores a $d\times d$ inverse curvature
estimate and two $d\times d$ covariance accumulators, so memory is $O(d^2)$; this is
inherited from the independent-data method we extend and is not something we improve. For
$d$ in the thousands this is the binding constraint, and a sketched or diagonal-plus-low-rank
curvature representation---as in \citet{kuang2026online,wang2026inference,godichonbaggioni2026complexity}---%
would be the natural remedy. Whether the block-gradient long-run covariance estimator
survives such a compression is open.

\paragraph{Geometric mixing is a real restriction, and the block length has to respect it.}
\Cref{thm:lrv} buys $b\asymp\log N$ from the geometric decay in
\Cref{lem:davydov}. Under algebraic mixing the flat-top bias is polynomial rather than
exponential and the admissible $b$ grows accordingly; long-memory streams are outside the
theory altogether. This is not merely formal: our air-quality stream has a residual lag-one
autocorrelation of $\RealAcfAirq$ and a median variance-inflation factor of $\RealKappaAirq$, and there the
block length needed for calibrated intervals is an order of magnitude larger than for the
traffic stream (\Cref{tab:real}). The free diagnostic
$r_j=|u_j^{\!\top}(\Sft-\Sbm)u_j|/(u_j^{\!\top}\Sft u_j)$ detects the problem---it estimates
the relative block-length bias of plain batch means---but it does not remove it. It could be
turned into an automatic rule by maintaining covariance accumulators at several block scales
$b,2b,4b,\dots$ at once (one extra $O(d^2)$ update per block per scale, all from the same
block gradients) and reporting from the smallest scale at which $r_j$ falls below a
tolerance; we did not implement this, and we do not claim an automatic rule for $b$.

\paragraph{$\Sft$ need not be positive semidefinite.}
This is intrinsic to flat-top and lugsail estimators \citep[Remark~5]{singh2025utility}.
Marginal intervals need only positive scalars $u_j^{\!\top}\Sft u_j$. The documented fallback to
$\Sbm$ never triggered in any simulation, but it does trigger on the real streams --- at a rate
of $\RealPsdFallback$ of coordinate-replicate pairs --- which is what one should expect, since
there the two designs are ill-conditioned and the lag-one correction is comparable to the
batch-means term. We report the rate rather than claiming it away, and a user who needs the full
matrix (for a joint test, say) must project it.

\paragraph{The curvature form is required.}
Everything rests on
$\nabla^2F(\theta)=\E[\alpha\,\Psi\Psi^{\!\top}]$, which holds for generalised linear
models, ridge and softmax regression, the geometric median and Gauss--Newton problems, but
not for a general objective. Without it there is no rank-one recursion and no
inversion-free update.

\paragraph{Warm-up is a real requirement, not an implementation detail.}
The $1/t$ schedule with a poor preconditioner leaves a slowly decaying component in the error.
On the autoregressive design the effect is now mild --- with a negligible warm-up, coverage of
nominal $95\%$ intervals is $\CovNoWarmup$ rather than $\CovBgsnArhom$ --- because the
step-length safeguard absorbs most of what a bad preconditioner does. On the regime-switching
design it is not mild: the relative variance is $\WarmMarkovFifty$ at $c_{\mathrm w}=50$ against
$\WarmMarkovTwoHundred$ at $c_{\mathrm w}=200$ (\Cref{tab:ablation}(c$'$)). We give a scaling rule
($c_{\mathrm w}d$ curvature updates, with $c_{\mathrm w}=200$ throughout) and show
insensitivity above it, but this is a tuning constant, and the asymptotic theory is silent
about it. What sets its size is not the number of observations but the number of effectively
distinct design points they span: on a persistent regime-switching stream, $50d$ consecutive
observations cover only about thirty regime visits, and the curvature estimate is still
rank-starved (\Cref{tab:ablation}(c$'$)). The warm-up also has a cost consequence we state
rather than bury: it performs $c_{\mathrm w}d$ rank-one updates at $O(d^2)$ each, so the total
time is $O(dN+c_{\mathrm w}d^3)$ and the advertised $O(dN)$ only bites once
$N\gtrsim c_{\mathrm w}d^2$. With $c_{\mathrm w}=200$ and $d=40$ that means
$N\gtrsim3\times10^5$, so at the $N=2\times10^5$ of \Cref{tab:cost} the warm-up dominates the
streaming pass for every $d\ge40$ we report, which is why that table separates the streaming
count from the total.

\paragraph{Blocking creates a stability condition, and our fix is not fully proved.}
The contraction condition $\gamma_t\opnorm{\Hb^{-1}\Hh_t}<2$ is a consequence of blocking the
gradient, so it is a cost of our construction and not of the method we build on. The step-length
safeguard in \eqref{eq:step} removes the failure it causes, and \Cref{prop:shift} shows that on
the event where the safeguard is eventually inactive it changes no asymptotic statement. What we
do \emph{not} prove is that this event has probability one: the implication we establish runs
from convergence to eventual inactivity, not back, because the safeguard factor depends on the
current block's gradient and so is not predictable, which breaks the martingale structure the
proofs use. \Cref{app:shift} says exactly what would close the gap --- an expanding-truncation
construction \citep{kushner2003stochastic,andrieu2005stability}, or a predictable variant that
caps using the previous block's amplification at an extra $O(d^2)$ per block --- and we did
neither. Instead we report the activation counts, which are between $2$ and $7$ per run and flat
in $N$ across every design and sample size in \Cref{tab:ablation,tab:hard}. Two further caveats.
The condition itself comes from the linearised recursion \eqref{eq:err}, so it is a heuristic for
non-quadratic objectives, exact only in the linear-model case. And $c_D$ is a tuning constant with a genuine failure mode:
\Cref{tab:ablation}(i) shows it is irrelevant over a factor of sixteen on the designs that do not
need the safeguard, but on the one that does, $c_D=4$ is too loose to prevent the blow-up. The
asymptotic theory says nothing about it, and we have no rule for choosing it beyond the value that
worked across our designs.

\paragraph{The central limit theorem is proved for fixed $b$, and we recommend growing $b$.}
\Cref{thm:clt} is unconditional at each fixed block length, but \Cref{thm:lrv} takes
$b\asymp\log N$, so the interval we recommend sits in the triangular-array regime where our proof
needs an extra hypothesis: that the finite random constants in \Cref{lem:series} are bounded in
probability uniformly in $N$ along the sequence $b_N$. We state that hypothesis in
\Cref{thm:clt} rather than hiding it, and we do not verify it. Closing the gap means redoing
\Cref{lem:series} with constants explicit in $b$; the naive route --- bounding
$\E\opnorm{H^{-1}-P_{t-1}}\norm{\eps_t}$ directly --- is the one the appendix shows loses a factor
$T^{1/2}$, which is why Gordin's decomposition is there in the first place. Two things make us
believe the composition is sound in practice rather than merely convenient: growing $b$ makes the
block scores less dependent, not more, so the mechanism the constants come from is working in our
favour; and coverage is flat in $b$ over $b\in[10,320]$ at $N=200{,}000$
(\Cref{tab:ablation}(b)), which is the statement the two theorems need to compose.

\paragraph{Rates and what we do not prove.}
We give an almost sure rate and a central limit theorem, not a Berry--Esseen bound; the
first-order Markovian literature has the latter for linear stochastic approximation
\citep{samsonov2025statistical}, and we do not match it. We attain the semiparametric
efficiency bound of \citet{li2023markovian} rather than establishing it. Our
$\widetilde O(N^{-1/2})$ rate for $\Sft$ is not comparable to the minimax rate
$\Theta(N^{-(1-a)/2})$ for Hessian-free covariance estimation from an iterate path
\citep{ni2026online}, because we use more information (block gradients and a maintained
curvature estimate) and because our step exponent is $a=1$, outside the range that bound is
stated for. Within the Hessian-free class the bound is close to attained
\citep{wei2026refining}, so the gap we report is about the information used, not about slack in
their analysis. We make no claim of optimality within our own information set: we do not know
the minimax rate for estimators that may use block gradients and a running curvature estimate.

\paragraph{Scope of the empirical evidence.}
All designs are convex $M$-estimation with $d\le160$ and $N\le4\times10^5$; there is no
deep-learning experiment and the theory would not support one. The dependence in the
simulated designs is stationary and geometrically decaying by construction. On the real
streams we do not know the truth, which is why we run two protocols; the semi-synthetic
protocol has an exact target but a synthetic (autoregressive) error process, and the fully
real protocol has a real error process but a target that is itself estimated---we correct
the nominal level for that exactly, but the correction assumes the segments are
approximately exchangeable, which non-stationarity would violate. Neither protocol alone
would be convincing; together they bracket the honest answer.

\paragraph{One thing we deliberately do not claim.}
Gapping improves the curvature estimate at matched cost (\Cref{prop:gap}, \RefGapFig{}),
but the curvature estimate enters the limit distribution only through a rate, so the
end-to-end effect on the point estimate and on coverage is second order and we could not
measure it reliably. The case for gapping is that it is free, deterministic, and provably
better than the randomised alternative---not that it changes the headline numbers.
